\problema{Meeting Rooms II}{253}{src/03-intervalos/253-meeting-rooms-ii.cpp}{M}

Nos dan un arreglo de intervalos de reunión \texttt{intervals[i] = [inicio,
fin]}. Queremos el número \emph{mínimo} de salas necesarias para que todas
las reuniones se puedan realizar, sabiendo que dos reuniones que se
traslapan en el tiempo no pueden compartir sala.

\paragraph{La idea, con un ejemplo de la vida real.} Imagina que eres el
encargado de reservar salas de juntas en una oficina, y te llegan todas las
solicitudes de reunión del día de golpe. En vez de intentar asignar salas
reunión por reunión, es más fácil pensar en dos listas separadas: ``a qué
horas empieza alguna reunión'' y ``a qué horas termina alguna reunión''.
Si recorres el día cronológicamente y llevas la cuenta de cuántas salas
están ocupadas en cada instante, el número máximo que alcanzas en algún
momento es la respuesta: necesitas al menos esa cantidad de salas para no
quedarte sin espacio.

\begin{center}
\ttfamily
\begin{tabular}{l}
0 -----------------------------30\\
\ \ \ \ 5----10\\
\ \ \ \ \ \ \ \ \ \ \ \ \ \ 15----20\\
\end{tabular}
\end{center}

\noindent Con \texttt{intervals = [[0,30],[5,10],[15,20]]}: \texttt{[5,10]}
y \texttt{[15,20]} no se traslapan entre sí, así que podrían compartir una
segunda sala; pero ambas se traslapan con \texttt{[0,30]}, que necesita su
propia sala durante todo ese tiempo. En total hacen falta \texttt{2} salas.

\paragraph{Cómo lo resuelve el código, paso a paso.}
\begin{enumerate}
  \item Separamos los \texttt{inicio}s de todas las reuniones en un
        arreglo, y los \texttt{fin}es en otro. A partir de aquí ya no
        importa a qué reunión pertenecía cada hora, solo el momento en que
        ocurre.
  \item Ordenamos ambos arreglos por separado, de menor a mayor.
  \item Recorremos el tiempo con dos punteros, uno para \texttt{inicio}s y
        otro para \texttt{fin}es (barrido / \emph{sweep line}). En cada
        paso comparamos el siguiente inicio pendiente contra el siguiente
        fin pendiente.
  \item Si el próximo inicio ocurre \emph{antes} de que la reunión más
        próxima a terminar acabe, necesitamos una sala adicional: sumamos
        uno al contador de salas en uso y avanzamos el puntero de inicios.
  \item Si no (el próximo fin ocurre antes o al mismo tiempo), una sala se
        libera: restamos uno al contador de salas en uso y avanzamos el
        puntero de fines.
  \item Llevamos registro del valor \emph{máximo} que alcanzó el contador de
        salas en uso durante todo el barrido: esa es la respuesta.
\end{enumerate}

\lstinputlisting{src/03-intervalos/253-meeting-rooms-ii.cpp}

\complejidad{$O(n \log n)$}{$O(n)$}

\paragraph{¿Qué significa esa complejidad?} Ordenar los dos arreglos de $n$
horas cada uno domina el costo, $O(n \log n)$; el barrido con dos punteros
después es lineal. En espacio guardamos dos arreglos de tamaño $n$ (los
inicios y los fines por separado).

\paragraph{Consejos para la entrevista (Amazon).} Un desempate que suele
pasar desapercibido: cuando un inicio y un fin ocurren \emph{exactamente al
mismo tiempo}, ¿libera la sala antes de que se necesite otra, o hace falta
una sala nueva de todos modos? La convención estándar (y la que produce la
respuesta correcta) es tratar los fines como si ocurrieran primero: una
reunión que termina a las 10 libera la sala justo a tiempo para otra que
empieza a las 10. Menciona también la alternativa con un min-heap de horas
de fin (equivalente en complejidad, más fácil de explicar en pizarrón), y
el caso límite de un arreglo vacío, que no necesita ninguna sala.
